%for a more compact document, add the option openany to avoid
%starting all chapters on odd numbered pages
%--------------------------------------------------------------------------
\documentclass[12pt]{cmuthesis}
%--------------------------------------------------------------------------
% This is a template for a CMU thesis.  It is 18 pages without any content :-)
% The source for this is pulled from a variety of sources and people.
% Here's a partial list of people who may or may have not contributed:
%
%        bnoble   = Brian Noble
%        caruana  = Rich Caruana
%        colohan  = Chris Colohan
%        jab      = Justin Boyan
%        josullvn = Joseph O'Sullivan
%        jrs      = Jonathan Shewchuk
%        kosak    = Corey Kosak
%        mjz      = Matt Zekauskas (mattz@cs)
%        pdinda   = Peter Dinda
%        pfr      = Patrick Riley
%        dkoes = David Koes (me)

% My main contribution is putting everything into a single class files and small
% template since I prefer this to some complicated sprawling directory tree with
% makefiles.

%--------------------------------------------------------------------------
% some useful packages
%--------------------------------------------------------------------------
\usepackage{times}
\usepackage{fullpage}
\usepackage{graphicx}
\usepackage{amsmath}
\usepackage[numbers,sort]{natbib}
\usepackage[backref,pageanchor=true,plainpages=false, pdfpagelabels, 
bookmarks,bookmarksnumbered,
%pdfborder=0 0 0,  %removes outlines around hyper links in online display
]{hyperref}
\usepackage{subfigure}
\usepackage{comment}   % package for block comment
\usepackage{graphicx}
\usepackage{hyperref}
\usepackage{url}
\usepackage{breakurl}
\usepackage{color}
\usepackage{xcolor}
\usepackage{listings}

% Approximately 1" margins, more space on binding side
%\usepackage[letterpaper,twoside,vscale=.8,hscale=.75,nomarginpar]{geometry}
%for general printing (not binding)

%--------------------------------------------------------------------------
\usepackage[letterpaper,twoside,vscale=.8,hscale=.75,nomarginpar,
hmarginratio=1:1]{geometry}
%--------------------------------------------------------------------------


\newcommand{\two}{\mspace{-2.0mu}}
\newcommand{\four}{\mspace{-4.0mu}}
\newcommand{\plus}{\mspace{-4.5mu}+\mspace{-3.5mu}}
\newcommand{\minus}{\mspace{-4.5mu}-\mspace{-3.5mu}}
\newcommand{\pp}{'\mspace{-2.0mu}'}

\newcommand{\xlb}[4]{#1\ifthenelse{\equal{#2}{0}}{}{_{\alpha #2}}
\mspace{-2.0mu}\genfrac{(}{)}{0pt}{1}
{\ifthenelse{\equal{#3}{0}}{0}{l #3}}
{\ifthenelse{\equal{#4}{0}}{0}{b #4}}}
%--------------------------------------------------------------------------
\newcommand{\xkv}[4]{#1\mspace{-5.0mu}\left(\mspace{-8.0mu}
\begin{smallmatrix}#2\four{}\four{}\mspace{-8.0mu}&
\pmb{\kappa}#3\\&\nu #4
\end{smallmatrix}\mspace{-5.0mu}\right)}

\newcommand{\evect}[6]{#1\mspace{-4.0mu}\left(\mspace{-8.0mu}
\begin{smallmatrix}#2\mspace{-8.0mu}&
\pmb{\kappa} #3 &b #5\\&\nu #4 &\alpha #6
\end{smallmatrix}\mspace{-5.0mu}\right)}

\newcommand{\varmat}[8]{\mspace{-5.0mu}\left(\mspace{-8.0mu}
\begin{smallmatrix}\ifthenelse{\equal{#3}{0}}
{\mspace{-8.0mu}&b_{#1}&b_{#2}\\&\alpha_{#1}&\alpha_{#2}}
{\ifthenelse{\equal{#7}{0}}{#1\mspace{-8.0mu}&\pmb{\kappa}#2#3
\mspace{-8.0mu}&\pmb{\kappa}#4#5\mspace{-8.0mu}&
\pmb{\kappa}#6\\&\nu#2&\nu#4&\nu#6} 
{#1\mspace{-8.0mu}&\pmb{\kappa}#2#3\mspace{-8.0mu}&\pmb{\kappa}#4#5
\mspace{-8.0mu}&\pmb{\kappa}#6#7\mspace{-8.0mu}&
\pmb{\kappa}#8\\&\nu#2&\nu#4&\nu#6&\nu#8}}
\end{smallmatrix}\mspace{-5.0mu}\right)}
%--------------------------------------------------------------------------
\newcommand{\EXP}[1]{\exp\mspace{-5.0mu}\left[#1\right]\mspace{-3.0mu}}

\newcommand{\tpp}[2]{\left(\mspace{-2.0mu}\xkv{\omega}{}{}{}#1
\xkv{\omega}{}{'}{'}#2\xkv{\omega}{}{\pp}{\pp}\mspace{-2.0mu}\right)}

\newcommand{\be} {\begin{eqnarray}}
\newcommand{\ee} {\end{eqnarray}}
\newcommand{\f}[2]{\ensuremath{\frac{\displaystyle{#1}}
{\displaystyle{#2}}}}
\newcommand{\lr}[1]{\langle{#1}\rangle}

\newcommand{\SUM}[2]{\ifthenelse{\equal{#1}{0}}
{\sum_{\alpha_{#2},b_{#2},l_{#2}}^{3,n,N}} 
{\ifthenelse{\equal{#1}{1}}{\sum_{\alpha_{#2},b_{#2}}^{3,n}}
{\sum_{\pmb{\kappa}#2,\nu#2}^{N,3n}}}}

\newcommand{\SUMprime}[2]{\ifthenelse{\equal{#1}{0}}
{\sum_{\alpha_{#2},b_{#2},l_{#2}}^{3,n,N}} 
{\ifthenelse{\equal{#1}{1}}{\sum_{\alpha_{#2},b_{#2}}^{3,n}}
{\sum_{\pmb{\kappa}^{'}#2,\nu#2}^{N,3n}}}}
%--------------------------------------------------------------------------
\newcommand{\SUMalpha}[2]{\ifthenelse{\equal{#1}{0}}
{\sum_{\alpha_{#2}}^{3}} 
{\ifthenelse{\equal{#1}{1}}{\sum_{\alpha_{#2},b_{#2}}^{3,n}}
{\sum_{\pmb{\kappa}#2,\nu#2}^{N,3n}}}}

\newcommand{\SUMalphap}[2]{\ifthenelse{\equal{#1}{0}}
{\sum_{\alpha'_{#2}}^{3}} 
{\ifthenelse{\equal{#1}{1}}{\sum_{\alpha'_{#2},b'_{#2}}^{3,n}}
{\sum_{\pmb{\kappa}#2,\nu#2}^{N,3n}}}}

\newcommand{\SUMb}[2]{\ifthenelse{\equal{#1}{0}}{\sum_{b_{#2}}^{n}} 
{\ifthenelse{\equal{#1}{1}}{\sum_{\alpha_{#2},b_{#2}}^{3,n}}
{\sum_{\pmb{\kappa}#2,\nu#2}^{N,3n}}}}

\newcommand{\SUMbp}[2]{\ifthenelse{\equal{#1}{0}}{\sum_{b'_{#2}}^{n}} 
{\ifthenelse{\equal{#1}{1}}{\sum_{\alpha'_{#2},b'_{#2}}^{3,n}}
{\sum_{\pmb{\kappa}#2,\nu#2}^{N,3n}}}}

\newcommand{\SUMl}[2]{\ifthenelse{\equal{#1}{0}}{\sum_{l_{#2}}^{N}} 
{\ifthenelse{\equal{#1}{1}}{\sum_{\alpha_{#2},b_{#2}}^{3,n}}
{\sum_{\pmb{\kappa}#2,\nu#2}^{N,3n}}}}

\newcommand{\SUMlp}[2]{\ifthenelse{\equal{#1}{0}}{\sum_{l'_{#2}}^{N}} 
{\ifthenelse{\equal{#1}{1}}{\sum_{\alpha'_{#2},b'_{#2}}^{3,n}}
{\sum_{\pmb{\kappa}#2,\nu#2}^{N,3n}}}}
%--------------------------------------------------------------------------
\newcommand{\abcdt}[5]{\mspace{-4.0mu}\left(\mspace{-8.0mu}
\begin{smallmatrix}&\ifthenelse{\equal{#1}{}}{a}{#1}&
\ifthenelse{\equal{#3}{}}{c}{#3}\\&
\ifthenelse{\equal{#2}{}}{b}{#2}&\ifthenelse{\equal{#4}{}}{d}{#4}
\end{smallmatrix}\mspace{-2.0mu};
\ifthenelse{\equal{#5}{}}{t}{#5}\right)}

\newcommand{\abcd}[4]{\mspace{-4.0mu}\left(\mspace{-8.0mu}
\begin{smallmatrix}&\ifthenelse{\equal{#1}{}}{a}{#1}&
\ifthenelse{\equal{#3}{}}{c}{#3}\\&
\ifthenelse{\equal{#2}{}}{b}{#2}&\ifthenelse{\equal{#4}{}}{d}{#4}
\end{smallmatrix}\mspace{-3.0mu}\right)}

\newcommand{\abt}[3]{\mspace{-4.0mu}\left(\mspace{-8.0mu}
\begin{smallmatrix}&\ifthenelse{\equal{#1}{}}{a}{#1} \\&
\ifthenelse{\equal{#2}{}}{b}{#2}
\end{smallmatrix}\mspace{-2.0mu};
\ifthenelse{\equal{#3}{}}{t}{#3}\right)}

\newcommand{\ab}[2]{\mspace{-4.0mu}\left(\mspace{-8.0mu}
\begin{smallmatrix}&\ifthenelse{\equal{#1}{}}{a}{#1} \\&
\ifthenelse{\equal{#2}{}}{b}{#2}
\end{smallmatrix}\mspace{-3.0mu}\right)}

\newcommand{\kvbat}{\mspace{-4.0mu}\left(\mspace{-8.0mu}
\begin{smallmatrix} &\pmb{\kappa} &b \\ &\nu &\alpha
\end{smallmatrix}\mspace{-2.0mu};t\right)}

\newcommand{\kvbatp}{\mspace{-4.0mu}\left(\mspace{-8.0mu}
\begin{smallmatrix} &\pmb{\kappa} &b' \\ &\nu &\alpha'
\end{smallmatrix}\mspace{-2.0mu};t\right)}
%--------------------------------------------------------------------------
\newcommand{\kvbaw}{\mspace{-4.0mu}\left(\mspace{-8.0mu}
\begin{smallmatrix} &\pmb{\kappa} &b \\ &\nu &\alpha
\end{smallmatrix}\mspace{-2.0mu};\omega\right)}

\newcommand{\kvbawp}{\mspace{-4.0mu}\left(\mspace{-8.0mu}
\begin{smallmatrix} &\pmb{\kappa} &b' \\ &\nu &\alpha'
\end{smallmatrix}\mspace{-2.0mu};\omega\right)}

\newcommand{\kvba}{\mspace{-4.0mu}\left(\mspace{-8.0mu}
\begin{smallmatrix} &\pmb{\kappa} &b \\ &\nu &\alpha
\end{smallmatrix}\mspace{-3.0mu}\right)}

\newcommand{\kvbap}{\mspace{-4.0mu}\left(\mspace{-8.0mu}
\begin{smallmatrix} &\pmb{\kappa} &b' \\ &\nu &\alpha'
\end{smallmatrix}\mspace{-3.0mu}\right)}

\newcommand{\kpvba}{\mspace{-4.0mu}\left(\mspace{-8.0mu}
\begin{smallmatrix} &\pmb{\kappa}^{'} &b \\ &\nu &\alpha
\end{smallmatrix}\mspace{-3.0mu}\right)}

\newcommand{\kva}{\mspace{-4.0mu}\left(\mspace{-8.0mu}
\begin{smallmatrix} &\pmb{\kappa} \\ &\nu &\alpha
\end{smallmatrix}\mspace{-3.0mu}\right)}

\newcommand{\kvap}{\mspace{-4.0mu}\left(\mspace{-8.0mu}
\begin{smallmatrix} &\pmb{\kappa} \\ &\nu &\alpha'
\end{smallmatrix}\mspace{-3.0mu}\right)}
%--------------------------------------------------------------------------
\newcommand{\kvb}{\mspace{-4.0mu}\left(\mspace{-8.0mu}
\begin{smallmatrix} &\pmb{\kappa} &b \\ &\nu 
\end{smallmatrix}\mspace{-3.0mu}\right)}

\newcommand{\kvbp}{\mspace{-4.0mu}\left(\mspace{-8.0mu}
\begin{smallmatrix} &\pmb{\kappa} &b' \\ &\nu 
\end{smallmatrix}\mspace{-3.0mu}\right)}

\newcommand{\kvt}{\mspace{-4.0mu}\left(\mspace{-8.0mu}
\begin{smallmatrix}&\pmb{\kappa} \\&\nu
\end{smallmatrix}\mspace{-2.0mu};t\right)}

\newcommand{\kpvt}{\mspace{-4.0mu}\left(\mspace{-8.0mu}
\begin{smallmatrix}&\pmb{\kappa}^{'} \\&\nu
\end{smallmatrix}\mspace{-2.0mu};t\right)}

\newcommand{\kvw}{\mspace{-4.0mu}\left(\mspace{-8.0mu}
\begin{smallmatrix}&\pmb{\kappa} \\&\nu
\end{smallmatrix}\mspace{-2.0mu};\omega\right)}

\newcommand{\kv}{\mspace{-4.0mu}\left(\mspace{-8.0mu}
\begin{smallmatrix}&\pmb{\kappa} \\&\nu
\end{smallmatrix}\mspace{-3.0mu}\right)}
%--------------------------------------------------------------------------
\newcommand{\lbt}{\mspace{-4.0mu}\left(\mspace{-8.0mu}
\begin{smallmatrix}&l \\&b\end{smallmatrix}\mspace{-2.0mu};t\right)}

\newcommand{\lbtp}{\mspace{-4.0mu}\left(\mspace{-8.0mu}
\begin{smallmatrix}&l' \\&b'\end{smallmatrix}\mspace{-2.0mu};t\right)}

\newcommand{\lt}{\mspace{-4.0mu}\left(\mspace{-8.0mu}
\begin{smallmatrix}&l\end{smallmatrix}\mspace{-2.0mu};t\right)}

\newcommand{\ltp}{\mspace{-4.0mu}\left(\mspace{-8.0mu}
\begin{smallmatrix}&l'\end{smallmatrix}\mspace{-2.0mu};t\right)}

\newcommand{\lb}{\mspace{-4.0mu}\left(\mspace{-8.0mu}
\begin{smallmatrix}&l \\&b\end{smallmatrix}\mspace{-3.0mu}\right)}

\newcommand{\lbp}{\mspace{-4.0mu}\left(\mspace{-8.0mu}
\begin{smallmatrix}&l' \\&b'\end{smallmatrix}\mspace{-3.0mu}\right)}
%--------------------------------------------------------------------------

% Provides a draft mark at the top of the document. 
\draftstamp{\today}{DRAFT}

\begin {document} 
\frontmatter

A) Signature Page – see Section 4
s
%initialize page style, so contents come out right (see bot) -mjz
\pagestyle{empty}

B) Title Page – see Section 4
\title{ %% {\it \huge Thesis Proposal}\\
{\bf Thermal Modeling of Disordered Materials}}
\author{Jason M. Larkin}
\date{August 2013}
\Year{2013}
\trnumber{}



\committee{
\href{http://www.cmu.edu/me/people/alan-j-h-mcgaughey.html}
{Alan J. H. McGaughey, Chair} \\
\href{http://www.cmu.edu/me/people/jonathan-a-malen.html}
{Jonathan Malen} \\
\href{http://www.ce.cmu.edu/people/faculty/maloney.html}
{Craig Maloney} \\
\href{http://www.cmu.edu/physics/people/faculty/widom.html}
{Michael Widom}
}

\support{}
\disclaimer{}

6.2 Preliminaries



C) Copyright Notice – should follow the title page, on a separate page, if statutory copyright in the dissertation has been or is to be claimed

%--------------------------------------------------------------------------
% copyright notice generated automatically from Year and author.
% permission added if \permission{} given.

\keywords{Phonon}

\maketitle

\begin{dedication}
For my dog
\end{dedication}

\pagestyle{plain} % for toc, was empty

%% Obviously, it's probably a good idea to break the various sections of your thesis
%% into different files and input them into this file...

\begin{abstract}
A short summary.
\end{abstract}

\begin{acknowledgments}
test
\end{acknowledgments}

D) Acknowledgments – must include financial support acknowledgment as discussed in Section 5

E) Abstract – should not exceed 350 words; see Section 2

F) Table of Contents – with page references

G) List of Tables – with titles and page references

H) List of Figures and Illustrations – with titles and page references.



\tableofcontents
\listoffigures
\listoftables
%--------------------------------------------------------------------------
\mainmatter

%% Double space document for easy review:
%\renewcommand{\baselinestretch}{1.66}\normalsize

% The other requirements Catherine has:
%
%  - avoid large margins.  She wants the thesis to use fewer pages, 
%    especially if it requires colour printing.
%
%  - The thesis should be formatted for double-sided printing.  This
%    means that all chapters, acknowledgements, table of contents, etc.
%    should start on odd numbered (right facing) pages.
%
%  - You need to use the department standard tech report title page.  I
%    have tried to ensure that the title page here conforms to this
%    standard.
%
%  - Use a nice serif font, such as Times Roman.  Sans serif looks bad.
%
% Other than that, just make it look good...

%--------------------------------------------------------------------------

I get the requirements here
\href{http://www.cit.cmu.edu/current_students/graduates/thesis_dissertation_policies.html}{http://www.cit.cmu.edu/current_students/graduates/thesis_dissertation_policies.html}

6. PREPARATION OF THE MANUSCRIPT: FORM

 

6.1 Title – Your thesis or dissertation will be a valuable source for other scholars only if it can be located easily. Modern retrieval systems use the words in the title—and sometimes a few other descriptive words—to locate your document. It is essential that the title be a meaningful description of the content of your dissertation. Avoid highly specialized terms to the extent possible, and use word substitutes for formulas, symbols, superscripts, Greek letters, etc.

 



 

6.3 Body of thesis or dissertation

A) Introduction

B) Main Body – with larger divisions and more important minor divisions indicated by suitable, consistent headings

C) Summary and Conclusions – highlighting the key findings and conclusions of the work presented. For engineering and science theses and dissertations, this section often also includes recommendations for follow-up research.

D) References – >see discussion in Section 6.9

E) Appendices – each appendix should have a title and be listed in the Table of Contents

 

6.4 Page Numbering

Each page in a thesis or dissertation should be assigned a number. The following plan of page numbering generally is accepted:

A) For Preliminaries – use small Roman numerals (i, ii, iii, iv, etc.). The numbering begins with ii; the title page counts as page i, but the number does not appear.

B) For Remainder of Thesis or Dissertation – including the text, illustrations, appendices, and bibliography, use Arabic numerals (1, 2, 3, 4, etc.). Each page must be numbered. Try to avoid use of letter suffixes such as 10a, 10b. The numbering begins with 1 and runs consecutively to the end of the dissertation.

C) For More Than One Volume – each volume should contain a title page duplicating the title page of the first volume. If the volumes are separate entities, identify them further as Volume I, II, etc. In any case, the numbering may follow consecutively from one volume to another, or begin with Arabic 1 at each new title page.

 

6.5 Footnotes

If footnotes are needed, they should be placed at the bottom of the page below a 1.5 inch underscore (starting at the left border). The first line of each footnote should be indented 0.5 inches and identified by a raised numeral corresponding to that used in the test. Footnotes should be numbered consecutively throughout each chapter.




\input{/home/jason/thesis/thesis/vc/vc}
\chapter{Vibrational Mean Free Paths and Thermal Conductivity Accumulation Functions for
Amorphous Materials}

test

\section{\label{mfp:Introduction}new section}



\chapter{Origin of the Exceptionally-Low Thermal Conductivity of Fullerene-Derived PCBM thin films}
There are two sections to this introductory chapter.


%--------------------------------------------------------------------------
\chapter{\label{Conclusion}Conclusion (WORK) (READ)}
%--------------------------------------------------------------------------

%--------------------------------------------------------------------------
\section{\label{Overview}Overview and Contributions}
%--------------------------------------------------------------------------

%--------------------------------------------------------------------------
\subsection{\label{Overview:SED} 
Molecular Dynamics-based Methods for Predicting Vibrational 
Lifetimes}
%--------------------------------------------------------------------------

In Chapter \ref{Chapter:SED}, two MD-based methods for predicting 
phonon properties and thermal conductivity were compared. 
The $\Phi$ method, which is the NMD method in the frequency-domain, 
was properly derived the starting with anharmonic lattice 
dynamics theory (see Appendix \ref{Appendix_A:Derivation}). 
The meaning of the proposed spectral method, $\Phi'$, 
was clarified and related it to the dynamic structure factor 
(see Appendix \ref{Appendix_B}). 
While 
the $\Phi'$ method does not accurately predict the mode lifetimes, 
the advantage of the $\Phi'$ versus the $\Phi$ method is that it does 
not require an eigenvalue solution for the mode eigenvectors.  
The dynamic structure factor, closely related to the $\Phi'$ method 
(see Appendix \ref{Appendix_B}), can predict frequency-dependent 
timescales from MD simulations for systemswith a larger number of 
atoms than those studied in this work using the NMD method 
(see Section \ref{Future:Timescales}). 

%--------------------------------------------------------------------------
\subsection{\label{Overview:VC}
Thermal Transport in Alloys and the High-scatter Limit}
%--------------------------------------------------------------------------

In Chapter \ref{chapter:vc}, thermal transport in two model alloys 
was investigated. The work provides several original 
insight into the physics of thermal 
transport in disordered lattices (i.e., isotopic solids and alloys). 
The first rigourous test of the virtual crystal (VC) approximation 
was presented. The VC-ALD technique has been used in a number of 
recent studies,
\cite{garg_role_2011,tian_phonon_2012,lindsay_thermal_2012,
li_thermal_2012,luckyanova_coherent_2012} 
but its limits had not been assessed until this study.  
The limits of the VC-ALD 
approach were determined using computationally-inexpensive empirical 
potentials and self-consistently treating the disorder explicitly 
and as a perturbation. 
The results indicate that while VC-ALD is generally 
an accurate method for materials 
whose thermal conductivity is dominated by low-frequency vibrational 
modes, care must be taken when modeling alloys with low thermal 
conductivities, where significant underprediction of thermal 
conductivity is likely. 

The high-scatter limit of thermal diffusivity,  
typically used in modeling amorphous materials, is directly relevant 
to the modeling of disordered lattices. 
Application of the VC 
approximation leads to vibrational mode diffusivities that are 
non-physical and the high-scatter limit provides a simple, 
physically-sound approach for correcting these predictions. 
The following calculations that were performed are novel additions 
to the literature:
\begin{itemize}
\item Using the VC-NMD method to model disordered 
lattices explicitly. This approach is limited in 
that the group velocity cannot be extracted so that thermal conductivity 
cannot be predicted.  The novel contribution is the use of 
NMD to predict the lifetimes of a disordered lattice using VC-NMD, 
where the normal modes of the Virtual Crystal (VC) are used as an 
approximation (see Section , Fig. and Appendix ).
\item  To model the disorder 
explicitly, the AF Theory calcualtions were perfomred 
on a disodered lattice. This theory has only previously been applied 
to amorphous phases.
\cite{feldman_thermal_1993,feldman_numerical_1999,
shenogin_,he_heat_2011}
The AF theory predictions showed that the lower-limit 
of diffusivity of high-frequency modes in a disordered lattice is the 
high-scatter limit, in contrast to the VC-ALD method, which incorrectly 
predicts that the limiting value is zero. Identification of this 
high-scatter limit of mode diffusivity was essential for identifying 
the breakdown in the VC-ALD method.  The high-scatter limit of 
diffusivity is usually assumed, without theoretical justification, in 
phenomological models for disordered and amorphous materials.
\cite{kittel_interpretation_1949,graebner_phonon_1986,
cahill_lattice_1988} 
This study gives self-consistent justification for its use.
\item Structure Factor of Disordered Lattice to Predict Group Velocities. 
The structure factor was calculated for modes in a model disordered 
lattice, which has previously been calculated for modes in amorphous 
materials.
\cite{biswas_vibrational_1988,feldman_thermal_1993,
allen_diffusons_1999,feldman_numerical_1999,
taraskin_determination_1999,taraskin_propagation_2000,
volz_molecular-dynamics_2000,
gotze_evolution_2000,horbach_high_2001,
martin-mayor_dynamical_2001,feldman_calculations_2002,
ciliberti_brillouin_2003,christie_vibrational_2007,
shintani_universal_2008,wyart_scaling_2010,
beltukov_ioffe-regel_2013,larkin_predicting_2013,
marruzzo_heterogeneous_2013} 
The structure factor predictions help us to 
understand that the VC-predicted group velocities are an 
underprediction of the representative velocity scale for mode 
diffusivities in the disordered lattice.  While previous studies 
have attempted to predict the group velocity of modes in disordered 
systems, there is no theoretical justification for the methods used.
\cite{duda_reducing_2011,donadio_atomistic_2009,
he_heat_2011,he_thermal_2011,he_morphology_2011,hori_phonon_2013} 
The structure factor provides a rigorous manner to estimate 
group velocities and is a significant contribution to understanding 
how to predict the correct velocity scale for mode diffusivities in 
disordered systems. 
\end{itemize}

By using all four methods discussed in 
Section \ref{S-SystemMode}, a self-consistent study of the VC approximation 
identified important connections between the modeling 
of disordered lattices and amorphous materials. 
By using three phases of LJ argon (perfect crystal, disordered lattice, 
and amorphous phase), the applicability 
of the different methods for predicting the thermal conductivity and 
mode-properties was demonstrated:
\begin{itemize}
\item MD-based GK method: suitable for modeling all 
three phases, but does not predict the mode properties.
\item Phonon based VC-ALD and VC-NMD: suitable for the perfect crystal 
and disordered lattices with the high-scatter limit correction.
\item The AF theory of diffusons: suitable for the high-frequency modes 
of the disordered lattice and all modes of the amorphous phase. 
We are unaware of any other study that uses all five of these methods 
self-consistently on the same material system. Our work provides clear 
guidelines for others on what tools are appropriate for different 
solid state systems.
Contribution to the Literature
Due to their low thermal conductivities, alloys are currently an active 
area of research, notably in the thermoelectric energy conversion field. 
The ability to predict alloy thermal conductivity is critical in 
narrowing down a large materials design space. Recent papers 
[e.g., PRL 106, 045901 (2011), PRL 109, 095901 (2012), 
PRB 85, 184303 (2012)] have used the VC-ALD method to make such 
predictions.
We believe this work will make an important contribution to the 
literature because the high-scatter limit adjustment is of interest 
to the study of low-thermal conductivity alloys. Thermoelectric energy 
generation materials, such as PbTe/Se alloys [21,22,49], maximize their 
efficiency by minimizing their thermal conductivity. The search for 
lower thermal conductivity alloys will require the modeling of even 
lower-thermal conductivity alloys, where the high-scatter limit we 
have proposed should be considered.
\end{itemize}

% The breakdown of the VC-ALD method has gone unnoticed in previous 
% computational work because these studies: 
% \begin{itemize}
% \item were limited to the 
% VC-ALD method because of computationally-expensive DFT calculations, 
% so that validation was not possible [12,18-27].
% \item focused on 
% materials where the thermal conductivity is dominated by 
% low-frequency vibration modes [12,18-27]
% \item did not always compare their 
% predictions with experimental measurements.\cite{tian_phonon_2012}
% \end{itemize}

% In our work, we provide a self-consistent study of thermal transport 
% in disordered lattices using a set of complementary computational tools 
% based in molecular dynamics simulations and lattice dynamics calculations. 
% The use of empirical potentials versus computationally-expensive DFT 
% calculations allowed us to perform the molecular dynamics simulations 
% that were necessary to observe the breakdown of the VC-ALD method. Our 
% study includes two test materials that demonstrate the applicability 
% and breakdown of the VC-ALD method. The conclusions are of general use 
% for the study of any disordered lattice. 

%--------------------------------------------------------------------------
\subsection{\label{Overview:MFP}
Mean Free Paths of Propagating Modes in Disordered Materials}
%--------------------------------------------------------------------------

A clear theoretical and modeling framework for amorphous materials was 
presented, which can form the basis for studying any disordered material. 
Drew on results and theory from the amorphous literature, which 
is extensive. 
Much work has been done by the physics community on vibrational spectra 
in amorphous solids, particularly on a-SiO$_2$. However, few of these 
studies investigated thermal transport properties. Only recently 
have thermal transport properties been studied.\cite{}

The mode lifetimes can be predicted using NMD.
If the model is large and low enough frequencies are accessed, the sound 
speed can be used. This can establish an upper bound on the mode 
diffusivities. 

The NMD lifetimes can be used with the AF theory diffusivities to predict 
effective diffuson MFPs and group velocities, which are shown in insets 
of Figs. \ref{F:vAF}(a) and \ref{F:vAF}(b). As noted in Section , the 
diffuson MFPs range between the system size $L$ and the lattice constant 
$a$ (see Section ). 
For diffusons, this effective diffuson velocities range from the 
sound speed to a fraction of it (about 0.3 for a-Si, see Fig. ). These 
diffuson velocities are much larger than 
those predicted in previous studies where the effective group velocities 
for some modes were approximately zero.\cite{}

\begin{itemize}
\item Identified the effects of metastability in amorphous materials 
on predicting lifetimes using the NMD method 
(see Section and Appendix ). 
Metastability is likely 
to affect the application of the NMD method in other ordered and 
disordered systems with weak bonding (see Section \ref{Future:LUC}). 
\item Identified differnces in the structural properties of 
of a-SiO$_2$ and a-Si that lead to a substantial difference in the 
propagating contribution to thermal conductivity in each. 
\item Predicted effective dispersion from the static structure factor (see 
Section ). The effective dispersion, when combined with diffusivity 
predictions from the AF theory and NMD method, are able to predict 
self-consistently the mode group velocities in a disorded system. 
Previous attempts to predict mode group velocities lack a theoretical basis. 
Presented a clear theoretical and computational framework 
for predicting the mode-level properties in amorphous materials. 
By predicting the effective dispersion using the static structrue 
factor, showed that the sound speed is the valid for 
low frequency modes. By using the NMD and AF-predicted 
diffusivities, demonstrated that the effective group velocity of 
diffuson modes is larger than predictions from recent studies.
(cite) 
\item Demonstrated that the NMD-predicted diffusivities are more 
reliable than those predicted by the AF theory at low frequencies 
(see Section ).
\item By comparing predictions from the models of a-Si using 
both $\omega^{-2}$ and $\omega^{-4}$ scalings, showed that 
further experimentation for a-Si thin films.

\end{itemize}


%--------------------------------------------------------------------------
\begin{figure}
\begin{center}
\centering
\includegraphics[scale=1.0]
{/home/jason/disorder/si/amor/m_af_si_normand_4096_Lambda_3.eps}
\end{center}
\caption{\label{F:vAF} 
MFPs of propagons and diffusons in a-SiO$_2$ and a-Si (see Section ).
}
\end{figure}
%--------------------------------------------------------------------------

\clearpage

%--------------------------------------------------------------------------
\subsection{\label{S:Predictive}
Predictive ability versus Computational Cost (WORK) (READ)}
%--------------------------------------------------------------------------

With the results from all studies...

now prepared to make a ranking of predictive capability. 

The GK method played an imporant role in verifying the mode properties 
prediced by all methods. In Section, the GK method provided a common 
comparison for the $\Phi$ and $\Phi'$ methods, which helps to confirm 
the disagreement between the two methods. In Section, the GK method 
provides a comparison for predictions from the VC-NMD and VC-ALD 
methods, which helps to identify the high-scatter limit of 
the diffuson mode diffusivities in disordered lattices.(cite) 
Finally, the GK method helps to confirm the scaling of the 
low-frequency contribution of the finite models of a-Si which 
is used to extrapolate to a bulk (infinte)-size value (Section ). 
The GK method will most likely be a vaulable modeling tool 
for future work. 

The VC-ALD method was shown to be limited to low frequency modes 
and best suited to high-thermal conductivity materials. High thermal 
conductivity (low-frequency dominated) materials are well-modeled 
by VC-ALD. VC-ALD is not well-suited to low thermal conductivity 
(full spectrum) materials. 
The AF theory models accurately the mid- and high-frequency 
modes in disordered materials, it does not handle the low-frequency 
modes for disordered lattices. The predictions for amorphous 
materials is also inconsistent. 
The VC-ALD method and AF theory can be supplimented by NMD calcualtions. 
The VC-NMD method is able to accurately predict the lifetimes 
of all vibrons in disordered lattices (see Section ). Propagating modes 
in a-Si can be identified definitively by NMD-predicted 
lifetimes (see Section ), while the AF theory does not show a 
definitive scaling (see Section ). 

With these findings, the theoretical methods (first discussed in 
Section ) are ranked in order of their predictive capabilites are in 
Table \ref{T-comparison-predictive}. The NMD, while the most 
computationally demanding of the four theoretical methods 
(Table \ref{T-comparison}), 
is ranked first in mode- and second in system-level predictive 
capability. The reasons for these rankings are:

\begin{itemize}
\item The NMD method is derived correctly from anharmonic lattice dynamics 
theroy (see Section ).. NMD predicts accurately the mode lifetimes and 
theraml conductivities compared to the $\Phi'$ method (see Section ).
\item The VC-NMD method accurately predicts the mode lifetimes for 
disordered lattices compared to the VC-ALD method. This leads to better 
agreement with the GK method, which is the most accurate system-level 
method.
\item The NMD method accurately predicts the low-frequncy lifetimes 
for a-Si, while the AF theory predictions have large fluctuations and 
depend on the broadening factor. The scaling from the NMD lifetimes is 
used to extrapolate a bulk thermal conductivity which is in good 
agreement with the system-level GK method (see Section ). 
\end{itemize}

The ALD and AF theory are considered to be equivalent at predicting 
mode-level and system-level properties because, for disordered lattices 
(alloys, see Section ), the VC-ALD method fails to accurately 
predict the mode lifetimes at high frequencies, 
while the AF theory is invalid at low frequencies 
(see Section ). Either 
method could be considered superior depending on whether the material 
being studied is low-frequency dominated or full-spectrum 
(see Section ). 

%--------------------------------------------------------------------------
\begin{center}
%\squeezetable
\begin{table}
\small
\caption{\label{T-comparison-predictive}Ranking of the predictive ability of 
theoretical techniques for mode-level and system-level thermal 
properties for disordered systems.}
%\begin{tabular}{llllll}
\begin{tabular}{p{0.9in}|p{0.9in}|p{0.9in}|p{0.9in}|p{0.9in}|p{0.9in}}
\hline\hline
System-level
&CH Theory 
&ALD 
&AF Theory 
&NMD 
&GK \\ 
\hline
Mode-level
&CH Theory 
&GK
&AF Theory
&ALD
&NMD \\ 
\hline\hline
\end{tabular}
\end{table}
\end{center}
%--------------------------------------------------------------------------

%--------------------------------------------------------------------------
\section{\label{Future}Future Work}
%--------------------------------------------------------------------------

%--------------------------------------------------------------------------
\subsection{\label{Future:LUC}Large Unit Cell Materials (WORK) (READ)}
%--------------------------------------------------------------------------

Large unit cell (LUC) materials are an important class of crystalline 
materials with a wide range of transport applications.(cite) 
Large unit cell crystals have an ordered (crystalline) structure, but 
the basic building block (unit cell) of the crystal has a large number 
of distinct atoms. 
They are effectively disordered 
over length scales on the order of the atomic spacing and their thermal 
conductivities can be as low as a glass.\cite{cahill_lower_1992} 
The key advantage of LUC materials is that they are still ordered 
from the standpoint of electrons, which results in good thermoelectric 
performance.
\cite{he_thermoelectric_2006,yang_effect_2006,wang_thermoelectric_2007} 

Some LUC 
materials, such as SiO$_2$-based zeolites, have been 
long-studied. Others, such as C$_60$(cite) or 
\href{http://en.wikipedia.org/wiki/Phenyl-C61-butyric_acid_methyl_ester}
{PCBM}, are currently being investigated for their thermal properties.
\cite{duda_exceptionally_2013} 
While experimental measurements of PCBM demonstrates that propagating 
modes contribute negligibly, the mechanism for its 
$\emph{exceptionally-low}$ is still not understood. Modeling could 
provide the necessary insights. 

From a modeling perspective, LUC materials pose a number of challenges,
theoretically and computationally, as compared to small unit cell 
(SUC) materials:

\begin{itemize}
\item LUC are crystalline, but are often organic or 
organic/inorganic hybrid materials. The structure of LUC 
materials is often poly- or quasi-crystalline, with less 
long-range order than SUC materials.(cite) 

\item Predicting model-level properties  
using ALD is challenging because the computational time 
scales as $b^4$ (see Section \ref{S-CompCost}). 

\item The presence of weak bonding in organic/inorganic 
materials can lead to metastability 
(see Appendix \ref{A:metastability}), which  
makes it challenging to perform the NMD method. 

\item MD simulations of LUC materials also face challenges. 
Based on the classical MD simulations of disordered 
materials studied in this work, systems sizes of 100-1000 atoms 
and simualtion times of up to 1 ns are expected to be 
necessary to accurately predict thermal conductivity. DFT calcualtions 
are too computationally expensive for such MD simulations.(cite)
Because many LUC materials have complex bonding environments, 
even the classical interatomic 
potentials are often computationally-expensive.(cite)
\end{itemize}

Based on the results in this work, there are several 
modeling strategies that can be used to study LUC materials:

\begin{itemize}
\item Identify signs of propagating modes from experimental 
measurements, if available.(cite) 
\item Based on the results for alloys (see Section ), the AF 
diffuson theory may have application for LUC materials, particularly 
at high frequencies and for 
those LUC materials which are only quasi-crystalline, 
such as PCBM.(cite)
\item The high-scatter limit for thermal conductivity (Eq. ) 
can be used to establigh a plasuible lower-bound for LUC 
materials.(cite) Similarly, the 
high-scatter limit for mode diffusivity can establish 
lower-bounds on the mode-level properties. 
\item Perform calcualtions using computationally-inexpensive 
classical interatomic potentials to asses if DFT calculations 
are necessary. 
\end{itemize}


%--------------------------------------------------------------------------
\begin{figure}
\begin{center}
\centering
\includegraphics[scale=0.4]
{/home/jason/thesis/thesis/duda_exceptionally_2013_fig4.eps}
\end{center}
\caption{\label{F:supercells} 
Range of thermal conductivity for crystalline and amorphous materials. 
Recent measurements on LUC materials show thermal conductivities 
much lower than bulk amorphous materials. 
Reproduced from Ref. \citenum{duda_exceptionally_2013}.
}
\end{figure}
%--------------------------------------------------------------------------

\clearpage

%--------------------------------------------------------------------------
\subsection{\label{Future:Timescales}
Timescales from Larger MD Simulations (WORK) (READ)}
%--------------------------------------------------------------------------

%--------------------------------------------------------------------------
\subsubsection{\label{Future:Timescales:NMD}
Exact Normal Modes}
%--------------------------------------------------------------------------

The NMD method used throughout this work is limited by the computational 
demands, which require a larger number ($\sim$ 100) of parallel processors 
to perform the analysis in a reasonable amount of time 
(less than 24 hours, see Appendix \ref{A:coding_lang:Execute:Scaling}). 
While the NMD method is trivially-parallelizable over the normal 
modes, the eigenvalue solution of the normal modes themselves 
is more computationally demanding. 
Standard routines for eigenvalue solutions of the dynamical matrix 
(see Section ) can 
calculate the exact normal modes for systems up to 8000 atoms 
in less than 24 hours using current computational resources.
\cite{gale_general_2003}  
These eigenvalue solution routines typically run on single 
processors. 
The eigenvalue solutions can be performed using the 
\href{http://www.mcs.anl.gov/petsc/}{Portable, Extensible 
Toolkit for Scientific Computation}(PETSc), which has routines for 
performing eigenvalue solutions in parallel. The PETSc pacakge has 
\href{http://wiki.python.org/moin/IntegratingPythonWithOtherLanguages}
{Python bindings} contained in the 
\href{https://code.google.com/p/petsc4py/}{petsc4py} 
pacakage, which allows for easy interface with existing MD packages
such as LAMMPS.\cite{plimpton_fast_1995}  
Here is a tutorial on 
\href{https://github.com/thehackerwithin/PyTrieste/wiki/F2Py}
{using the f2py package}, and a general tutorial on 
\href{http://docs.scipy.org/doc/numpy/user/c-info.python-as-glue.html}
{using Python as a ``glue'' language}. 

The controlling Python script would work as follows...

\begin{itemize}
\item The GULP package could be 
compiled as a Python module using the 
\href{https://code.google.com/p/f2py/}{f2py} package. 
\item The dynamical 
matrix could then be transferred to the parallel eigenvalue solution 
routines contained in PETSc using petsc4py.
\item The atomic trajectories are generated using the LAMMPS 
package with the liblammps Python module.
\item The 
\end{itemize}

A previous studied investigated the vibrational normal modes 
which has been used in a previous study on LJ argon to compute a 
(cite) 
small subset of the normal modes for a system of 32,000 atoms
\cite{mazzacurati_low-frequency_1996} 
using the 
\href{http://en.wikipedia.org/wiki/Lanczos_algorithm}
{Lanczos algorithm}.\cite{golub_matrix_2012} 
The same 32,000 system was subsquently studied using MD simulations 
and the dynamic structure factor,\cite{ruocco_relaxation_2000} 
which was discussed in Section .

While parallel eigenvalue solution can increase the system sizes 
accesible with NMD, the method is ultimately limited by the 
poor scaling of the run time and memory requirements (see Section ). 

%--------------------------------------------------------------------------
\subsubsection{\label{Future:Timescales:DSF}
Dynamic Structure Factor}
%--------------------------------------------------------------------------

MD simulations are computationally efficient.  Systems sizes of nearly 
$10^6$ atoms have been studied in this work, which were bulk systems 
with equal simulation side lengths in all three spatial dimensions. 
The dynamic structure factor, described in Section \ref{Appendix_B}, 
can predict vibrational timescales and does not 
need the eigenvectors of the exact normal modes to perform the mapping 
of the atomic trajectories. This method, 
combined with appropriately shaped supercells, could probe the 
timescales of vibrational modes up to wavelengths between 
24 and 100 nm using similar computational resources to those 
used in this work (see Section ). 
This presents an oppurtunity to compare with recent experimental 
measurements spectral linewidths in the frequency range below 
1 THz, which have been recently reported for a-SiO$_2$ but are 
lacking for a-Si.\cite{hondongwa_ultrasonic_2011} 
The current correlation function, closely related to the dynamic 
structure factor,\cite{horbach_high_2001} can even be used to study 
the spectral character of motions in a fluid.\cite{boom_molecular_1980}  

%--------------------------------------------------------------------------
\subsection{\label{Future:DFT}DFT and Thermal Modeling}
%--------------------------------------------------------------------------

Try MD and AF using better scaling DFT packages such as CP2K and BigDFT. 

BigDFT is linear scaling, but not until a large system size is reached. 

%--------------------------------------------------------------------------
\subsection{\label{Future:Package}
Comprehensive Package for Thermal Transport Calculations (WORK) (READ)}
%--------------------------------------------------------------------------

Four different methods were used in this work. No single package 
exists to perform the necessary calcaultions for all these methods. 


There exist numerous packages for performing calculations necessary for 
predicting thermal transport properties. No one package can perform all 
calculations necessary, particularly both the 
mode-level (bottom-up) and system-level (top-down) thermal transport 
properties.   

\begin{center}
\begingroup
%\squeezetable
\begin{table}
\caption{\label{T:available_codes}
Available packages
}
%\begin{ruledtabular}
\begin{tabular}{llllll}
\hline
Package & Language & Capabilities \\
\hline
GULP & Fortran & mode-level (harmonic) \\
LAMMPS & C++ & system-level, mode-level (limited) \\
phonopy & Python & mode-level (harmonic) \\
\end{tabular}
%\end{ruledtabular}
\end{table}
\endgroup
\end{center}

LAMMPS, for example, contains most methods for predicting the 
system-level thermal conductivity. A package to predict the 
mode-level properties is needed.  Ideally, the mode-level properties 
could be predicted alongside the system-level calculations, as is 
necessary to perform the NMD and spectral techniques described in 
Sections and . 

The Python language is an ideal environment for ''gluing`` together 
the availble codes, and extending their functionality in dynamic ways. 

LAMMPS as a model of social interaction, mailing list. 


num-gulp-citations(year) = [1996 1997 1998 1999 2000 2001 2002 2003 2004 2005 2006 2007 2008 2009 2010 2011 2012 2013] 

num-gulp-citations-year = [1 1 1 1 15 65 59 68 88 77 92 98 114 116 139 128 139 79*(12/9)] 



\appendix
\input{/home/jason/thesis/thesis/appendix/appendix_nmd}
%--------------------------------------------------------------------------
\chapter{\label{Appendix_HPC}
Research using High Performance Computing}
%--------------------------------------------------------------------------

%--------------------------------------------------------------------------
\section{\label{A:Comp_Env}Setting Up Computing Environment}
%--------------------------------------------------------------------------

% %%--------------------------------------------------------------------------
% %\subsection{Suggested Reading}
% %%--------------------------------------------------------------------------
% 
% %The online encyclopedia \href{http://www.wikipedia.org/}{wikipedia} 
% %has become an excellent resource for technical knowledge. 
% 
% %\href{http://en.wikipedia.org/wiki/Statistical_mechanics}
% %{Statistical Mechanics}/
% %\href{http://en.wikipedia.org/wiki/Condensed_matter}
% %{Condensed Matter}
% %\cite{ashcroft_solid_1976,mcquarrie_statistical_2000}
% 
% %\href{http://en.wikipedia.org/wiki/Phonon}
% %{Lattice Dynamics} with focus on the classical, quantum 
% %and thermodynamic properties of phonons.\cite{peierls_,ziman,
% %wallace_thermodynamics_1976,
% %srivastava_1990,dove_lattice_1993} 
% 
% %\href{http://en.wikipedia.org/wiki/Introduction_to_quantum_mechanics}{Introduction to quantum mechanics and quantum chemistry}
% %\cite{griffiths_introduction_1995} 
% %and more 
% %advanced topics on 
% %\href{https://en.wikipedia.org/wiki/Electron_configuration}
% %{Electron Strucutre} and 
% %\href{http://en.wikipedia.org/wiki/Density_functional_theory}
% %{Density Functional Theory}.
% %\cite{martin_electronic_2004} 
% 
% %Analytical methods for 
% %\href{http://en.wikipedia.org/wiki/Ordinary_differential_equations}
% %{ordinary} and 
% %\href{http://en.wikipedia.org/wiki/Partial_differential_equation}
% %{partial differential equations}, 
% %\href{http://en.wikipedia.org/wiki/Fourier_analysis}
% %{fourier} and 
% %\href{http://en.wikipedia.org/wiki/Statistics}
% %{statistical analysis},\cite{mcquarrie_mathematical_2003} 
% %and 
% %\href{http://en.wikipedia.org/wiki/Numerical_analysis}
% %{numerical analysis}.\cite{moin_fundamentals_2010} For example, 
% %the LJ MD code discussed in Section uses a  
% %\href{http://en.wikipedia.org/wiki/Verlet_integration}
% %{verlet integration method}. 

%--------------------------------------------------------------------------
\subsection{\label{A:Comp_Env:OS}Hardware and Operating System Choice}
%--------------------------------------------------------------------------

The choice of computing hardware determines which operating system(s) 
can be used. The three 
main choices for operating system are 
\href{http://en.wikipedia.org/wiki/Microsoft_Windows}{Windows}, 
\href{http://en.wikipedia.org/wiki/Mac_OS}{Apple OS}, and 
\href{http://en.wikipedia.org/wiki/Linux}{Linux}. Each 
operating system has limitations depending on the hardware it 
operates on.   
For example, Apple OS is primarily used on 
\href{http://en.wikipedia.org/wiki/Macintosh}{Apple hardware}.

The Linux operating system is based on unix, which is the  
\href{http://www.youtube.com/watch?v=7XTHdcmjenI}
{world's most widely used software}. 
Linux runs on systems ranging from 
\href{http://www.pcworld.com/article/250899/mobile_linux_its_not_just_android_anymore.html}
{cell phones} to the world's 
\href{http://www.zdnet.com/20-great-years-of-linux-and-supercomputers-7000018681/}
{largest supercomputers}. For personal computing use, the most widely 
used Linux version is \href{http://www.ubuntu.com/}{Ubuntu}.
It is well-documented with a large community-based discussion 
system. Apple OS is an adequate substitute as it is 
\href{https://en.wikipedia.org/wiki/Unix-like}{unix-like}.

There are many options for installing Ubuntu and the instructions 
can be found at 
\href{http://www.ubuntu.com/}{ubuntu.com}.
Ubuntu is certified on 
\href{http://www.ubuntu.com/certification/}
{PC's from many computer companies}. 
One company, \href{https://www.system76.com/}{system 76}, 
builds computers with Ubuntu pre-installed. I used the 
\href{https://www.system76.com/laptops/model/panp9}
{Pangolin Performance} for over a year of my Ph.D.  However, 
I recommend using the most-portable (i.e., lightest) and 
longest-battery life notebook available such as the 
\href{https://help.ubuntu.com/community/SamsungSeries9}
{Samsung Series 9}  
or \href{https://help.ubuntu.com/community/MacBookAir}{Macbook Air}. 
You will typically be using this notebook to access large computing 
clusters remotely, 
so there is a benefit of having portability and long 
battery life over computational power.  

%--------------------------------------------------------------------------
\subsection{\label{A:Comp_Env:Term}
Terminal, Commands, and Environment Variables}
%--------------------------------------------------------------------------

These instructions work best for the Ubuntu operating system, and will work 
well for other versions of Linux. Systems commands are executed by the 
system 
\href{https://help.ubuntu.com/community/UsingTheTerminal}{terminal}. 
The following Linux system commands have been used extensively throughout 
this work: sudo, cd, pwd, ls, export, echo, grep, sed, vi, cat, 
scp, ssh, which, nohup. 
There are many more Linux commands that are useful. More information can 
be obtained simply by google searching, i.e. 
\href{https://www.google.com/search?q=linux+sed}
{google: linux sed} 

\href{http://en.wikipedia.org/wiki/Environment_variable}{Environment variables}, 
such as 
\href{http://en.wikipedia.org/wiki/PATH_(variable)}{PATH}, 
can store the paths to installed programs and files. 
Here is how to determine what PATH is set to:
\begin{lstlisting}
$ echo $PATH
/opt/Datathief:/opt/VESTA-x86_64:/home/jason/phonopy/phonopy-1.1/bin:...
\end{lstlisting}
Permanent changes to environment variables can be made in the user's 
\href{https://www.gnu.org/software/bash/manual/html_node/Bash-Startup-Files.html}
{.bashrc file}, which is typically located in the /home/user/ directory. 
The .bashrc file is a 
\href{http://www.ghacks.net/2009/04/16/linux-tips-view-hidden-files/}{hidden file} 
noted by its name starting with ".". 
Here is an example 
\href{https://gist.github.com/kparrish/6064111}{.bashrc} from 
\href{https://github.com/kparrish}{Kevin Parrish}. 
This file modifies environment variables such as 
\href{https://help.ubuntu.com/community/EnvironmentVariables}{PATH}  
when a bash terminal session is launched. Changes to environment 
variables are made when a new terminal session is started. 

Shell scripts are relatively low-level sets of system commands 
which can manipulate environment variables 
and the Linux operating system. 
\href{http://linuxcommand.org/wss0010.php}{Here is a simple tutorial} 
on writing shell scripts. 
\href{http://docs.python.org/2/library/os.html}
{Running Linux system commands in Python} can be an effective way of 
generating and manipulating many files with one script. Python 
is a more robust language than lower-level shell scripting. 
\href{https://github.com/ntpl/ntpy/blob/master/examples/thesis/run.lammps.py}
{Here is an example} that generates a series of shell scripts used to 
submit jobs on cluster resources (see next section). The Python script 
finds dummy variable names (SHELL$\_$NAME and PROC$\_$NAME) and replaces them 
with a different string. 

%--------------------------------------------------------------------------
\subsection{\label{A:Comp_Env:Remote}Working on Remote Resources}
%--------------------------------------------------------------------------

At some point during research you will need to execute code on remote 
resources that are typically large computing clusters ($> 100$ central 
processing units). 
You will be provided with a terminal session similar to the session 
provided by Ubuntu with most of the same system commands. 
While I recommend 
\href{https://filezilla-project.org/}{Filezilla} 
for handling the transfer of data and files, 
the functionality of Filezilla is contained in several shell 
commands: ssh and scp. 
Files can be transferred using the following command
\begin{lstlisting}
$ scp -r user@gilgamesh.cheme.cmu.edu:/home/user/directory/ ./
\end{lstlisting}
which will place the ``directory'' and its contents into the pwd 
(./) of the local terminal session. 

There are many variants of the operating systems used for remote 
computing clusters, but the differences are usually small. 
During my work I regularly used the 
\href{http://gilgamesh.cheme.cmu.edu/doc/gilgamesh.html}
{gilgamesh} computing cluster maintained by 
\href{https://github.com/jkitchin}{John Kitchin}. 
\href{http://gilgamesh.cheme.cmu.edu/doc/gilgamesh.html}
{Gilgamesh's documentation} is a good resource for learning how to run 
calculations on a computing cluster.

%--------------------------------------------------------------------------
\section{\label{A:Comp_Env:Programs}
Installing, Writing and Executing Programs}
%--------------------------------------------------------------------------

%--------------------------------------------------------------------------
\subsection{\label{A:Comp_Env:Install}Installing Available Packages}
%--------------------------------------------------------------------------

Before writing any of your own code, it is best to use any useful 
code that may already exist. 
Most Linux operating systems have automatic software installation. 
For Ubuntu, program  management is achieved using the 
\href{https://help.ubuntu.com/community/AptGet/Howto}{apt-get} 
system command:

\begin{lstlisting}
$ sudo apt-get install gfortran
[sudo] password for jason: 
Reading package lists... Done
...
\end{lstlisting}
To check that the package has been installed, use:
\begin{lstlisting}
$ which gfortran
/usr/bin/gfortran
\end{lstlisting}
This shows that gfortran program is installed in the /usr/bin/ folder, 
a common location where automatically installed programs are located. 
Because the 
folder /usr/bin is in the PATH environment variable, 
the program gfortran is always available no matter where you are in 
the system terminal. 

To get more information about the gfortran, use:
\begin{lstlisting}
$ ls -l /usr/bin/gfortran
lrwxrwxrwx 1 root root 12 Apr 22 03:44 /usr/bin/gfortran -> gfortran-4.7
\end{lstlisting}
where we see that /usr/bin/gfrotran "points" to gfortran-4.7. 
This is called a 
\href{http://en.wikipedia.org/wiki/Symbolic_link}{symbolic link}.

%--------------------------------------------------------------------------
\subsection{\label{A:Comp_Env:Avail}Available Packages for 
Thermal Transport Modeling}
%--------------------------------------------------------------------------

Available programs represent opportunities to perform research quickly 
and easily. I suggest reading their documentation carefully and trying 
the tutorials and examples, which are typically computationally 
inexpensive. 

Most of the existing packages listed below cannot be installed using 
Ubuntu's apt-get command, although 
\href{http://lammps.sandia.gov/download.html#ubuntu}
{this feature does exist for LAMMPS} and more packages are likely 
to be added in the future. It is important to read the installation 
documentation carefully. It is also helpful to contact any system 
Linux/Unix system administration staff or other researchers who are 
experienced with installing packages. These people can usually be 
found at campus or local computing centers, or in your own 
research group! 

For a first-time user, I recommend trying to install the LAMMPS 
package as it is one of the easiest and most versatile. It has 
standard installation files for systems ranging from your 
personal computer up to a massively-parallel supercomputer. 
Installation instructions can be 
\href{http://lammps.sandia.gov/doc/Section_start.html#start_2}{found here}. 

Here is a partial list of the packages I have found useful 
in this work:

\begin{itemize}
\item \href{http://lammps.sandia.gov/}{LAMMPS}, including the particularly 
useful 
\href{http://lammps.sandia.gov/threads/threads.html}{mailing list} 
and
\href{http://lammps.sandia.gov/doc/Section_python.html}
{python interface}: parallel molecular dynamics using empirical interatomic 
potentials.
\item \href{http://www.stfc.ac.uk/CSE/randd/ccg/software/DL_POLY/25526.aspx}
{DL$\_$POLY}: parallel molecular dynamics using empirical interatomic 
potentials.
\item \href{http://projects.ivec.org/gulp/}{GULP}: harmonic lattice 
dynamics using empirical interatomic potentials. 
\item \href{http://phonopy.sourceforge.net/}{phonopy}: harmonic 
lattice dynamics with interface for $\emph{ab initio}$ package. 
\item \href{http://www.homepages.ucl.ac.uk/~ucfbdxa/phon/}{PHON}: 
harmonic lattice dynamics with interface for $\emph{ab initio}$ package.
\item \href{http://www.abinit.org/}{ABINIT}: $\emph{ab initio}$ package. 
\item \href{http://www.quantum-espresso.org/}{Quantum Espresso}: 
$\emph{ab initio}$ package. 
\item \href{https://www.vasp.at/}{VASP}: $\emph{ab initio}$ package. 
\item \href{http://icmab.cat/leem/siesta/}{SIESTA}: 
$\emph{ab initio}$ package. 
\item \href{http://www.msg.ameslab.gov/gamess/}{GAMESS}: 
$\emph{ab initio}$ package.
\item \href{http://www.cp2k.org/}{CP2K}: 
$\emph{ab initio}$ package.
\item \href{http://www.icams.de/content/departments/ams/madsen/boltztrap.html}{BOLTZTRAP} 
\item \href{http://www.ks.uiuc.edu/Research/vmd/}{VMD} 
(with 
\href{https://sites.google.com/site/akohlmey/software/topotools}
{topotools}): structure and visualization. 
\item \href{http://jp-minerals.org/vesta/en/}{VESTA}: structure and 
visualization. 
\end{itemize}

I recommend trying this 
\href{https://gist.github.com/kparrish/6064159}{install.sh} created by 
\href{http://www.github.com/kparrish}{Kevin Parrish} which installs 
many programs and packages, including LAMMPS for parallel computing 
with 
\href{http://www.open-mpi.org/}{OpenMPI}.

%--------------------------------------------------------------------------
\subsection{\label{A:Comp_Env:Writing}Writing Programs: 
Compiled versus Interpreted}
%--------------------------------------------------------------------------

If no package exists to perform the calculation you are interested in, 
there is a good chance you are doing something interesting! The next 
step is to implement the algorithm for the calculation you want 
to perform. There are two choices for doing this: 
\begin{itemize}
\item Write your own code. 
\item Build off an existing package. 
\end{itemize}

As outlined in the previous 
section, there are many pre-existing packages with many functionalities. 
The many packages that can perform nanoscale transport calculations use 
\href{http://en.wikipedia.org/wiki/Compiled_language}{compiled} and 
\href{http://en.wikipedia.org/wiki/Interpreted_language}{interpreted}  
languages, often with a mixture of the two. 
The most commonly used compiled languages are 
\href{https://en.wikipedia.org/wiki/C\%2B\%2B}{C/C++} 
(e.g., Linux, LAMMPS) 
and 
\href{http://en.wikipedia.org/wiki/Fortran}{Fortran}
(e.g., GULP, Quantum Espresso, VASP). 
Here is a good discussion on 
\href{http://stackoverflow.com/questions/13078736/fortran-vs-c-does-fortran-still-hold-any-advantage-in-numerical-analysis-thes}
{the strengths of C++ versus Fortran}. Here are excellent tutorials 
of programming in  
\href{http://www.youtube.com/watch?v=XFQ9dw3CyDo&list=PL1D10C030FDCE7CE0}
{C++}
and 
\href{http://www.youtube.com/watch?v=YRTEOFMUTzw&list=PL6A8E21D2E86A0155}
{Fortran}. Learning these languages is essential to adding to existing 
packages. 

Interpreted languages allow for fast code development and debugging. 
The two interpreted languages you are likely to use are 
\href{http://en.wikipedia.org/wiki/MATLAB}{Matlab}  
and 
\href{http://en.wikipedia.org/wiki/Python_(programming_language)}{Python}. 
The key to maximizing the potential of interpreted languages is by 
using the built-in ``vector'' functions and operations provided by the 
\href{http://www.mathworks.com/help/matlab/matlab_prog/vectorization.html}
{Matlab}  
and 
\href{http://faculty.washington.edu/rjl/uwamath583s11/sphinx/notes/html/python_vect.html}
{Python}  
programming libraries. 
Matlab has an excellent built-in guide and google searches will typically 
yield useful results. An open-source substitute for Matlab is 
\href{http://www.gnu.org/software/octave/}
{Octave}, which is capable to running most Matlab scripts.

The following three case studies are performed on my local 
\href{http://www.samsung.com/us/computer/series-9-notebooks}
{Samsung series 9 notebook}. The 
specifications of the computer can be accessed by using:
\begin{lstlisting}
$ cat /proc/cpuinfo 
processor	: 0
vendor_id	: GenuineIntel
cpu family	: 6
model		: 58
model name	: Intel(R) Core(TM) i5-3317U CPU @ 1.70GHz
stepping	: 9
microcode	: 0x17
cpu MHz		: 782.000
cache size	: 3072 KB
physical id	: 0
siblings	: 4
core id		: 0
cpu cores	: 2
apicid		: 0
initial apicid	: 0
\end{lstlisting}

%--------------------------------------------------------------------------
\subsubsection{\label{A:coding_lang:case1}Case-study (a): 
Optimizing Compiled Codes}
%--------------------------------------------------------------------------

The first case study is a 
\href{https://github.com/jasonlarkin/disorder/tree/master/md_serial}
{single C code to perform molecular dynamics 
on LJ argon}. The code uses simple subroutines and operates using 
a single (serial) processor. The 
\href{http://www.cplusplus.com/doc/tutorial/arrays/}
{arrays in this code are built statically}. 
Arrays can be created 
\href{http://www.cplusplus.com/doc/tutorial/dynamic/}{dynamically}, 
which allows for the input of systems with varying number of 
atoms. Additionally, 
\href{http://www.cplusplus.com/reference/vector/vector/}{vectors} 
can be created and destroyed dynamically and have some advantages 
over arrays.  

The code is compiled using the 
\href{http://www.gnu.org/}{GNU project's} 
C++ compiler, \href{http://www.cprogramming.com/g++.html}{g++}:
\begin{lstlisting}
$ g++ ArgonMD.cpp -o ArgonMD_O_g++
\end{lstlisting}
The code can be run and the output can be directed using a useful 
shell operation called 
\href{http://www.linfo.org/pipes.html}{piping}, 
which is demonstrated below:
\begin{lstlisting}
$ ArgonMD_O_g++ > ./ArgonMD_O/output.txt
\end{lstlisting}
The code can be compiled using 
\href{http://gcc.gnu.org/onlinedocs/gcc/Optimize-Options.html#Optimize-Options}
{optimization flags, such as -O and -O3}:
\begin{lstlisting}
$ g++ -O3 ArgonMD.cpp -o ArgonMD_O3_g++
\end{lstlisting}
which greatly decreases the run time of this particular code. 
The total run time 
can displayed by using the shell command 
\href{http://linux.about.com/od/commands/l/blcmdl1_grep.htm}{grep},  
which shows for no optimization:
\begin{lstlisting}
$ grep -A 1 "Total Time" ./ArgonMD/output.txt
Total Time: 38.42 (s)
\end{lstlisting}
for -O optimization:
\begin{lstlisting}
$ grep -A 1 "Total Time" ./ArgonMD_O/output.txt
Total Time: 18.06 (s)
\end{lstlisting}
and for -O3 optimization:
\begin{lstlisting}
$ grep -A 1 "Total Time" ./ArgonMD_O3/output.txt
Total Time: 11.74 (s)
\end{lstlisting}
A useful shell command is 
\href{http://www.tuxfiles.org/linuxhelp/vimcheat.html}{vi}, a 
shell-based text editor, which can display the output within 
a shell terminal:
\begin{lstlisting}
$ vi ./ArgonMD/output.txt
0       1       5.61358 2.01864 5.26596
...
\end{lstlisting}

To compare with the above results, the LAMMPS code is compiled in serial 
using no , -O, and -O3 optimizations.  
The results are contained in folders beginning with ``lmp''.  A useful 
shell function is 
\href{http://www.tuxfiles.org/linuxhelp/tabtrick.html}{tab-twice}, 
where tapping the tab key twice 
will display files and folders which begin with the same 
characters:
\begin{lstlisting}
$ grep -A 1 "Loop time" ./lmp
lmp.in.lj~     lmp_serial/    lmp_serial_O/  lmp_serial_O3/
\end{lstlisting}
The tab key can also be used for  
\href{http://en.wikipedia.org/wiki/Command-line_completion}
{command-line completion} to complete directory and file names. The run 
time for no optimization is:
\begin{lstlisting}
$ grep -A 1 "Loop time" ./lmp_serial/log.lammps
Loop time of 0.718188 on 1 procs for 1000 steps with 256 atoms
--
Loop time of 3.4892 on 1 procs for 5000 steps with 256 atoms
\end{lstlisting}
for -O optimization:
\begin{lstlisting}
$ grep -A 1 "Loop time" ./lmp_serial_O/log.lammps
Loop time of 0.19842 on 1 procs for 1000 steps with 256 atoms
--
Loop time of 0.917175 on 1 procs for 5000 steps with 256 atoms
\end{lstlisting}
and for -O3 optimization:
\begin{lstlisting}
$ grep -A 1 "Loop time" ./lmp_serial_O3/log.lammps
Loop time of 0.164066 on 1 procs for 1000 steps with 256 atoms
--
Loop time of 0.786311 on 1 procs for 5000 steps with 256 atoms
\end{lstlisting}

There is a decrease in run time with increasing optimization 
for LAMMPS and my code. However, for every 
optimization the LAMMPS code is approximately an order of 
magnitude faster.  This result is 
due to a several factors, the most important being 
the implementation of 
\href{http://en.wikipedia.org/wiki/Neighbor_list}{neighbor lists} 
as discussed in the 
\href{http://lammps.sandia.gov/doc/neighbor.html}
{LAMMPS documentation}. The interested reader is encouraged to 
investigate the LAMMPS code further for useful C++ coding 
practices. 

%--------------------------------------------------------------------------
\subsubsection{\label{A:coding_lang:case2}
Case-study (b): Computational Cost of Interpreted Code}
%--------------------------------------------------------------------------

With interpreted languages, the traditional programming practice of using 
loops (for/do/while, etc) will slow the code execution.  To illustrate 
this behavior, the case study compares codes to perform harmonic lattice 
dynamics calculations. 
\href{https://github.com/jasonlarkin/disorder/blob/master/matlab/m_lj_ld_prim.m}
{One code is written in Matlab} using mixture of loops 
and vectorized functions, which slows the code execution. The other 
code is written in Fortran (i.e., the GULP package). 

From within a Matlab session in a terminal, the Matlab code can be 
executed as follows:
\begin{lstlisting}
tic
[freq,eigVsorted,velocity] = m_lj_ld_prim(1.5636*[1 1 1],[0.5 0.0 0.0])
toc
Elapsed time is 0.047081 seconds.
\end{lstlisting}
for the 
\href{https://github.com/jasonlarkin/disorder/blob/master/matlab/m_lj_ld_prim.m}
{primitive}, and
\begin{lstlisting}
tic
[freq,eigVsorted,velocity] = m_lj_ld_conv(1.5636*[1 1 1],[0.5 0.0 0.0])
toc
Elapsed time is 0.260508 seconds.
\end{lstlisting}
for the 
\href{https://github.com/jasonlarkin/disorder/blob/master/matlab/m_lj_ld_conv.m}
{conventional} unit cells. Note that the increased run time is not 
proportional to the increased number of atoms in the unit cell (i.e., 
the calculation does not scale linearly). 
Also note that the two codes both use for loops. This usage slows down code 
execution because Matlab is an interpreted language. 

Now we run the same calculation using the Fortran-based GULP package, for 
the 
\href{https://github.com/jasonlarkin/disorder/blob/master/matlab/gulp_disp_lj_prim.gin}
{primitive}:
\begin{lstlisting}
$ grep -A 6 "Timing analysis " gulp_disp_lj_prim.gout
  Timing analysis for GULP :

---------------------------------------------------------------------------
  Task / Subroutine                                          Time (Seconds)
---------------------------------------------------------------------------
---------------------------------------------------------------------------
  Total CPU time                                                  0.0030
\end{lstlisting}
and 
\href{https://github.com/jasonlarkin/disorder/blob/master/matlab/gulp_disp_lj_conv.gin}
{conventional} unit cells:
\begin{lstlisting}
$ grep -A 6 "Timing analysis " gulp_disp_lj_conv.gout
  Timing analysis for GULP :
---------------------------------------------------------------------------
  Task / Subroutine                                          Time (Seconds)
---------------------------------------------------------------------------
---------------------------------------------------------------------------
  Total CPU time                                                  0.0160
\end{lstlisting}

First, the scaling of the calculation is non-linear. Second, the timing for 
both calculations is over an order of magnitude less than the Matlab versions. 
Matlab code can be optimized by using the built in vector operations and 
functions, which is demonstrated in the next section.

%--------------------------------------------------------------------------
\subsubsection{\label{A:coding_lang:case3}
Case-study (c): Optimizing Interpreted Code}
%--------------------------------------------------------------------------

This case study demonstrates that Matlab (or other interpreted languages) 
can be optimized using the built-in vector functions and operations. 
\href{https://github.com/jasonlarkin/disorder/blob/master/matlab/m_af_template.m}
{This Matlab code} performs the AF theory calculations 
on a LJ amorphous phase of 256 atoms. The code can be run from within 
a Matlab terminal using:
\begin{lstlisting}
>> m_af_template
ans =
    1.5075 (s)
\end{lstlisting}
which calls several subroutines, most notably 
\href{https://github.com/jasonlarkin/disorder/blob/master/matlab/m_af_lj.m}{m$\_$af$\_$lj.m}, 
which uses all vectorized functions. 

Comparing this result with the 
\href{https://github.com/jasonlarkin/disorder/blob/master/matlab/gulp/x0_256.gin}
{same calculation in GULP}:
\begin{lstlisting}
$ gulp x0_256
$ grep -A 10 "Timing analysis " x0_256.gout
  Timing analysis for GULP :
---------------------------------------------------------------------------
  Task / Subroutine                                          Time (Seconds)
---------------------------------------------------------------------------
  Calculation of real space energy and derivatives                0.0520
  Calculation of phonons                                         18.5720
  Calculation of scattering                                       0.0000
---------------------------------------------------------------------------
  Total CPU time                                                 19.3920
---------------------------------------------------------------------------
\end{lstlisting}
The relevant subroutine in GULP to compare with the Matlab version is 
thermalconductivity.f90. 
The GULP timing is considerably larger than the Matlab code.  This is 
due to the use of the default unoptimized 
\href{http://www.netlib.org/blas/}{BLAS} and 
\href{http://www.netlib.org/lapack/}{LAPACK} libraries 
for linear algebra that were used to compile GULP. During the installation, 
Matlab optimizes its linear algebra library, resulting 
in improved performance. 
\nomenclature[A]{BLAS}{Basic Linear Algebra Subprograms} 
\nomenclature[A]{LAPACK}{Linear Algebra PACkage} 

% %--------------------------------------------------------------------------
% \subsection{\label{A:coding_lang:Execute}Executing Programs}
% %--------------------------------------------------------------------------
% 
% %--------------------------------------------------------------------------
% \subsubsection{\label{A:coding_lang:Execute:Local}
% Executing Locally: Rapid Development}
% %--------------------------------------------------------------------------
% 
% %--------------------------------------------------------------------------
% \subsubsection{\label{A:coding_lang:Execute:Remote}
% Executing Remotely: Portable Batch Systems}
% %--------------------------------------------------------------------------

%--------------------------------------------------------------------------
\subsection{\label{A:coding_lang:Execute:Scaling}
Scaling Calculations}
%--------------------------------------------------------------------------

The majority of the methods 
used in this work scale poorly with the number of atoms, $N_a^\alpha$, 
with $\alpha >1$. 
The scaling cost for eigenvalue solution is $N_a^3$. 
The cost of performing this calculation for the largest system size,  
$cost_{large}$,  
and every successive system which is half the size of the largest 
is given by the 
\href{http://en.wikipedia.org/wiki/Geometric_series}{geometric series} 
with common ratio $r=1/8$
\begin{equation}\label{EQ:cost_total}
\begin{split}
cost_{total} = cost_{large}\frac{1}{1-r} = 1.143N_{a,large}. 
\end{split}
\end{equation}
It costs a minimal amount ($ \approx 14\%$) to study every system that 
is smaller 
than the largest system considered. Even a linear scaling   
has (with $r=1/2$) a total cost $cost_{total} = 2cost_{large}$. 
As such, I recommend starting calculations with 
the smallest system of interest (i.e., the 
\href{http://meta.tex.stackexchange.com/questions/3300/minimum-working-example-mwe}
{Minimum Working Example} method). 
Journal article draft manuscripts can be 
developed quickly by performing computationally-inexpensive 
calculations first, documenting the results, and then iterating to more 
computationally-expensive calculations. 
This scheme for performing calculations can follow these time-scales 
for calculation costs:
\begin{itemize}
\item 1 second: exploratory work, debugging.  
\item 1 minute/1 hour: testing the scaling of algorithms, precision  
\item 1 day: publication-quality calculations
\item 1 week: publication-quality calculations, but should be limited 
to when necessary. 
\end{itemize}
I have performed 
countless calculations costing around one second, and very few which cost 
more than one week. Publication quality results will typically 
cost between one hour and one week.
\nomenclature[V]{$cost_{total}$}{total computational cost} 

%--------------------------------------------------------------------------


%%--------------------------------------------------------------------------
%\subsection{Importance of Open-Source Code}
%%--------------------------------------------------------------------------

%LAMMPS Mailing List
%http://lammps.sandia.gov/threads/threads.html

%search: 

%lennard (jones)	45
%silicon		125
%nanotube	153
%ger


%Execellent discussion of ensembles and newtonian dynamics
%http://lammps.sandia.gov/threads/msg13979.htmlmanium	2
%silica		79
%carbon		181
%hydrogen	99
%fullerene(C60)	7(12)	
%Pb (lead)	2
%water		541
%copper(cu)	25
%gold		81
%diamond		43

%phase transition	21
%green kubo	18


%Execellent discussion of ensembles and newtonian dynamics
%http://lammps.sandia.gov/threads/msg13979.html

%video of python development tree

%\href{http://www.youtube.com/watch?v=3poNeQHUKrs}
%{google+ history with gource}

%\href{http://www.youtube.com/watch?v=cNBtDstOTmA}
%{python history with gource}

%grand-daddy of them all, 
%\href{http://www.youtube.com/watch?v=pOSqctHH9vY}
%{linux kernel history with gource}

%It is important that results become more open-source.  It's important 
%that our communication becomes open-source. It's important that the 
%entire numerical process be carried out open-source. 


%%--------------------------------------------------------------------------
%\subsection{Redundancy}
%%--------------------------------------------------------------------------
%
%There are 3 different codes (shiomi/esfarjani, broido, ankit) for 
%performing the same basic calculations.
%
%For LD, there is GULP, phonopy, SIESTA, etc.
%
%%--------------------------------------------------------------------------
%\subsection{Code Development}
%%--------------------------------------------------------------------------
%
%took ankit 10 months to re-develop esfarjani code.  
%
%code development time can be drastically reduced using pre-existing code. 
%codes written in a modular fashion can be added to easily.
%
%%--------------------------------------------------------------------------
%\subsection{Experiment Pre-dating Simulation}
%%--------------------------------------------------------------------------
%The ideal goal of simulation is to pre-date experiment.  This has not 
%been achieved yet.  See Fig.
%
%findings for:
%CNT/Graphene
%Si
%Thermoeletric (LUC, alloys, SL, etc)
%Perovskites
%PCM
%
%%--------------------------------------------------------------------------
%\subsection{Experiment Pre-dating Simulation}
%%--------------------------------------------------------------------------
%
%ntpl most cited paper 1-5:
%maradudin, ladd, dove, ziman, 

%--------------------------------------------------------------------------
\section{Preparing Journal Articles and Thesis}
%--------------------------------------------------------------------------

The job of the student is to prepare, submit, and publish 
peer-reviewed journal articles. There are many journals suitable for 
nanoscale transport topics. All of them accept 
\href{http://www.latex-project.org/}{Latex} prepared manuscripts. 
I recommend the Latex editor 
\href{http://kile.sourceforge.net/}{Kile}, while the simple 
\href{https://projects.gnome.org/gedit/}{gedit} works well and comes 
pre-installed with Ubuntu. Here is a 
\href{http://mally.stanford.edu/~sr/computing/latex-example.html}
{simple latex example} and how to generate a 
\href{http://tex.stackexchange.com/questions/1596/how-to-compile-a-latex-document}
{portable document format (PDF)} from the latex document.  

I recommend editing a written research document at least once every week. 
The exchange of such a document with the advisor can be facilitated using 
\href{https://www.dropbox.com/}{Dropbox} or 
\href{https://github.com/}{Github}. 
Github offers advantages over Dropbox, such as version control and a 
\href{https://github.com/ntpl/ntpy/wiki}{built-in online wiki}. 
Such a research document can turn into journal articles, 
which then turn into the thesis.  

To maintain the reference library I 
recommend 
\href{http://www.zotero.org/}{zotero}.  Here is an example 
\href{https://github.com/jasonlarkin/disorder/blob/master/paper/vc/jap/ntpl-032013.bib}
{reference.bib} 
file which is exported automatically by zotero. 
The references are generated from the latex document using 
\href{http://www.bibtex.org/Using/}{bibtex}, which compiles the contents of 
the reference.bib.

\href{https://github.com/jasonlarkin/disorder/blob/master/paper/vc/jap/jap_vc_jl_060413.tex}
{Here is an example} of the latex file 
used to create an article published from this work. 
This file uses 
\href{http://publish.aps.org/revtex}{revtex}, 
which is an article class used by $\emph{Physical Review}$, 
$\emph{Journal of Applied Physics}$, and others.
\href{https://github.com/robsimmons/cmu-thesis}
{Here is a Carnegie Mellon thesis template}, 
which was used to create 
\href{https://github.com/jasonlarkin/thesis/tree/master/thesis}
{this thesis document}.

% \href{http://web.science.mq.edu.au/~rdale/resources/writingnotes/latexstruct.html}
% {Article on structuring large documents.}
% %--------------------------------------------------------------------------
% \subsection{Producing Figures}
% %--------------------------------------------------------------------------
% 
% Python has a plotting module, \href{http://matplotlib.org/}{matplotlib}, 
% which has \href{http://matplotlib.org/examples/index.html}{many examples}. 
% Here is a \href{https://gist.github.com/kparrish/6101681}{simple example 
% demonstrating how to generate and save data, load that data, and plot.} 
% 
% 
% \href{https://github.com/ntpl/ntpy/tree/master/examples/thesis}
% {Here is a simple example demonstrating how to }

%--------------------------------------------------------------------------
\section{Technical Advice}
%--------------------------------------------------------------------------

In addition to your advisor and close mentors, 
I recommend communicating with experts in the field as much as 
possible (without being annoying).  How often to 
communicate depends on the situation. 
It is best to let the expert dictate the pace of the conversation. 
Here is a list of experts I used as resources for this work. They 
will typically answer emails within 24-48 hours:

\begin{itemize}
\item \href{http://www.engineering.pitt.edu/WissamAlSaidi/}
{Wissam Al-Saidi}
\item \href{http://www2.mpip-mainz.mpg.de/~donadio/tnt/People.html}
{Davide Donadio}
\item \href{http://johncduda.com/}
{John Duda}
\item \href{http://mech.rutgers.edu/content/keivan-esfarjani}
{Keivan Esfarjani}
\item \href{mailto:joseph.feldman.ctr@nrl.navy.mil}
{Joseph Feldman}
\item \href{http://chemistry.curtin.edu.au/people/academic.cfm/J.Gale}
{Julian Gale}, \href{http://projects.ivec.org/gulp/news.html}
{(GULP author)}
\item \href{http://quasiamore.mit.edu/pmwiki.php?n=Main.JivteshGarg}
{Jivtesh Garg}
\item \href{https://github.com/ankitjainmeiitk}
{Ankit Jain}
\item \href{https://github.com/jkitchin}
{John Kitchin}
\item Axel Kolhmeyer via the 
\href{http://lammps.sandia.gov/threads/threads.html}{LAMMPS mailing list} 
and \href{akohlmey@gmail.com}{email}
\item \href{http://jasonlarkin.github.io}{Jason Larkin}
\item \href{http://www.cmu.edu/me/malen/Lab_Website/People.html}
{Jon Malen}
\item \href{http://www.ce.cmu.edu/people/faculty/maloney.html}
{Craig Maloney}
\item \href{http://ntpl.me.cmu.edu/people.html}
{Alan McGaughey}
\item \href{http://www.pmc.umontreal.ca/~mousseau/site_an/}
{Normand Mousseau}
\item Steve Plimpton via the 
\href{http://lammps.sandia.gov/threads/threads.html}{LAMMPS mailing list} 
\item \href{http://www.chem.polimi.it/people/faculty/guido-raos/}
{Guido Raos}
\item \href{http://utexas.academia.edu/DanSellan}
{Dan Sellan}
\item \href{http://www.phonon.t.u-tokyo.ac.jp/}
{Junchiro Shiomi}
\item \href{http://atztogo.users.sourceforge.net/}{Atz Togo}, creator of 
\href{http://phonopy.sourceforge.net/}{phonopy}
\end{itemize}





\backmatter

%\renewcommand{\baselinestretch}{1.0}\normalsize

% By default \bibsection is \chapter*, but we really want this to show
% up in the table of contents and pdf bookmarks.
\renewcommand{\bibsection}{\chapter{\bibname}}
%\newcommand{\bibpreamble}{This text goes between the ``Bibliography''
%  header and the actual list of references}
\bibliographystyle{unsrt}
\bibliography{/home/jason/Dropbox/ntpl-paper/ntpl-060313} %your bib file
%--------------------------------------------------------------------------
\end{document}
